\section{Machine Learning Foundations for Heuristic Analysis}
\label{sec:ml_foundations}
\hspace{\parindent}
Although the main focus of this dissertation lies in centralized multi-robot path planning and
coordination, some of the methods used to design and evaluate the proposed heuristics rely on
basic machine learning concepts. In particular, clustering methods are relevant for analyzing the
spatial distribution of robots, while regression models can be used to study and estimate the
computational cost of heuristic functions as a function of relevant problem parameters.

This section introduces the machine learning concepts required to support these components.
The goal is not to provide an exhaustive review of machine learning methods, but rather to
establish the background necessary to understand the clustering-based heuristics and the runtime
modelling procedures used in the experimental evaluation.

\subsection{Supervised and Unsupervised Learning}
\label{subsec:supervised_unsupervised_learning}
\hspace{\parindent}
Machine learning methods are commonly divided into supervised and unsupervised learning
approaches \cite{theodoridis2020machine}. In supervised learning, a model is trained from examples composed
of input variables and known target outputs. The objective is to learn a mapping that can predict
the target output for new, unseen inputs. Regression is a typical supervised learning task, where
the target variable is continuous. In the context of this work, regression can be used to model the
execution time of heuristic functions based on input parameters such as the number of robots or
graph size.

In contrast, unsupervised learning methods operate on data without predefined target labels.
Their objective is to identify structure, patterns, or groups within the data. Clustering is one of the
most common unsupervised learning tasks, aiming to group samples that are similar according to
a given distance or similarity criterion \cite{theodoridis2020machine}.

\subsection{Clustering Methods}
\label{subsec:clustering_methods}
\hspace{\parindent}
Clustering methods aim to partition a set of data points into groups, called clusters, such that
points within the same cluster are more similar to each other than to points belonging to different
clusters. The notion of similarity is usually defined through a distance metric, such as the Euclidean
distance, although the specific choice depends on the representation of the data and the objective
of the analysis \cite{theodoridis2020machine}.

In multi-robot coordination, clustering can be used to characterize the spatial distribution of
robots in the environment. For example, robots located close to each other may compete for the
same spatial resources, create local congestion, or increase the probability of future interactions.
For this reason, clustering provides a useful tool for designing prioritization heuristics that account
not only for individual robot properties, but also for the local structure of the fleet.

In this context, two representative clustering approaches are particularly relevant: K-means and
DBSCAN. K-means is considered because it is a classical centroid-based clustering method and
provides a simple reference approach for grouping data points according to spatial proximity
\cite{macqueen1967kmeans,lloyd1982least}. However, it requires the number of clusters to be
defined in advance, which may be difficult in multi-robot planning scenarios where the number of
spatial robot groups changes across planning instances. For this reason, the elbow method is also
introduced as a common heuristic for selecting the number of clusters in K-means
\cite{thorndike1953family}.

DBSCAN is then presented as an alternative density-based method that does not require the
number of clusters to be specified beforehand and can explicitly identify isolated points as noise
\cite{ester1996dbscan}. 

\subsubsection{K-means}
\label{subsubsec:kmeans}
\hspace{\parindent}
K-means is a classical centroid-based clustering algorithm that partitions a set of data points
into \(k\) clusters, where \(k\) is defined in advance \cite{macqueen1967kmeans,lloyd1982least}.
Due to its simplicity, interpretability, and computational efficiency, it is one of the most widely
used clustering methods and is often adopted as a reference approach when discussing spatial
grouping problems. Given a set of observations \(X = \{x_1, x_2, \ldots, x_n\}\), the objective of
K-means is to assign each point to one of \(k\) clusters so as to minimize the sum of squared
distances between each point and the centroid of the cluster to which it belongs.

This objective can be written as

\begin{equation}
    \min_{C_1,\ldots,C_k}
    \sum_{j=1}^{k}
    \sum_{x_i \in C_j}
    \|x_i - \mu_j\|^2
\end{equation}
where \(C_j\) denotes the set of points assigned to cluster \(j\), and \(\mu_j\) is the centroid of
that cluster. The algorithm typically alternates between assigning each point to the nearest centroid
and updating each centroid as the mean of the points assigned to it. This process is repeated until
the assignments stabilize or until a maximum number of iterations is reached.

K-means is simple and computationally efficient, which makes it widely used in practice.
However, it has important limitations. First, the number of clusters \(k\) must be specified before
running the algorithm. Second, all points are forced to belong to one of the clusters, even if they
are isolated or do not naturally belong to any group. Third, because clusters are represented by
centroids, K-means is better suited to compact and approximately spherical clusters, and may
perform poorly when the underlying groups have irregular shapes.

These limitations are relevant in multi-robot path planning scenarios. The number of spatial
groups formed by robots is not usually known in advance and may vary between planning requests.
Moreover, isolated robots should not necessarily be forced into a cluster, since doing so may
artificially influence the priority score of unrelated robots.

\subsubsection{Elbow Method}
\label{subsubsec:elbow_method}
\hspace{\parindent}
Since K-means requires the number of clusters \(k\) to be specified beforehand, a common
procedure for selecting \(k\) is the elbow method \cite{thorndike1953family}. This method consists
of running K-means for several values of \(k\) and computing the corresponding within-cluster sum
of squares. As \(k\) increases, this value generally decreases, since more clusters allow points to be
represented by nearer centroids.

Let \(W(k)\) denote the within-cluster sum of squares obtained for a given value of \(k\):

\begin{equation}
    W(k) =
    \sum_{j=1}^{k}
    \sum_{x_i \in C_j}
    \|x_i - \mu_j\|^2.
\end{equation}

The elbow method selects a value of \(k\) corresponding to the point where increasing the number
of clusters produces only a limited additional reduction in \(W(k)\). This point is referred to as
the ``elbow'' of the curve.

Although the elbow method is intuitive and easy to apply, it is also subjective. In some datasets,
the curve does not exhibit a clear elbow, making the choice of \(k\) ambiguous. For this reason,
the elbow method is useful as an exploratory tool, but it does not remove the need to define or
estimate the number of clusters. This is one of the motivations for considering density-based
clustering methods in the context of spatial robot grouping.

\subsubsection{DBSCAN}
\label{subsubsec:dbscan}
\hspace{\parindent}
DBSCAN, or Density-Based Spatial Clustering of Applications with Noise, is a density-based
clustering algorithm that identifies clusters as dense regions of points separated by regions of lower
density \cite{ester1996dbscan}. Unlike K-means, DBSCAN does not require the number of clusters
to be specified in advance. Instead, it relies on two parameters: a distance threshold \(\varepsilon\)
and a minimum number of points \(MinPts\).

Given a point \(x_i\), its \(\varepsilon\)-neighbourhood is defined as

\begin{equation}
    N_{\varepsilon}(x_i)
    =
    \{x_j \in X : \|x_j - x_i\| \leq \varepsilon\}.
\end{equation}

A point is classified as a core point if its \(\varepsilon\)-neighbourhood contains at least \(MinPts\)
points. A point that is not a core point but lies within the \(\varepsilon\)-neighbourhood of a core
point is classified as a border point. Points that are neither core points nor border points are
classified as noise.

Clusters are then formed by connecting density-reachable points. Intuitively, a cluster grows
from a core point by recursively including neighbouring points that satisfy the density condition.
This allows DBSCAN to identify clusters with arbitrary shapes and to explicitly represent isolated
points as noise.

These properties are particularly useful for robot prioritization. In a planning request, robots may
form dense local groups in some regions of the environment while other robots remain isolated.
Since DBSCAN does not require a predefined number of clusters and does not force every robot
to belong to a cluster, it is well suited for detecting spatial concentrations of robots without
introducing artificial group assignments.

\subsubsection{Comparison between K-means and DBSCAN}
\label{subsubsec:kmeans_dbscan_comparison}
\hspace{\parindent}
K-means and DBSCAN represent two different approaches to clustering
\cite{macqueen1967kmeans,ester1996dbscan}. K-means is a centroid-based method that assumes
the number of clusters is known beforehand and assigns every point to one of the \(k\) clusters. It
is computationally efficient and works well when clusters are compact and clearly separated.
However, its dependence on a predefined \(k\) and its inability to represent noise make it less
suitable when the number of groups is unknown or when isolated points are meaningful.

DBSCAN, on the other hand, defines clusters according to local point density. It can discover
an unknown number of clusters, handle irregular cluster shapes, and identify isolated points as
noise. These characteristics make it more appropriate for the spatial analysis of robot configurations,
where the number and shape of robot groups may change from one planning instance to another.

In the context of this dissertation, this distinction is important because clustering is not used as
a general data analysis tool, but as part of a prioritization mechanism. The objective is to identify
spatially dense groups of robots that may create local congestion or competition for shared
resources. Since DBSCAN naturally captures such dense regions without requiring the number of
clusters in advance, it provides a better fit for this purpose than K-means.

\subsection{Linear Regression for Runtime Modelling}
\label{subsec:linear_regression_runtime}
\hspace{\parindent}
Regression methods are supervised learning techniques used to model the relationship between
a set of explanatory variables and a continuous target variable \cite{theodoridis2020machine,montgomery2021linear}.
Linear regression is one of the simplest and most interpretable regression models. It assumes that
the target variable can be approximated as a linear combination of the input variables.

Given a set of input features \(x_1, x_2, \ldots, x_p\), a multiple linear regression model estimates
a target variable \(y\) as

\begin{equation}
    \hat{y}
    =
    \beta_0
    +
    \sum_{j=1}^{p}
    \beta_j x_j,
\end{equation}

where \(\hat{y}\) is the predicted value, \(\beta_0\) is the intercept, and \(\beta_j\) are the model
coefficients associated with each input feature.

The coefficients are commonly estimated by minimizing the residual sum of squares between
the observed values \(y_i\) and the predicted values \(\hat{y}_i\):

\begin{equation}
    RSS
    =
    \sum_{i=1}^{n}
    (y_i - \hat{y}_i)^2.
\end{equation}

In this dissertation, linear regression can be used as an exploratory model for analyzing the
execution time of heuristic functions. The target variable corresponds to the measured runtime,
while the explanatory variables may include parameters such as the number of robots, the number
of vertices and edges in the roadmap, the value of \(k\) in \(k\)-hop neighbourhood heuristics, the
number of candidate priority orderings, or other scenario-dependent properties.

The main advantage of linear regression in this context is interpretability. Each coefficient
indicates how the predicted runtime changes with the corresponding input variable, assuming the
remaining variables are fixed. This makes the model useful not only for prediction, but also for
understanding which parameters have the strongest influence on computational cost.

However, linear regression also has limitations. It assumes a linear relation between the input
features and the target variable, which may not fully capture the behaviour of algorithms whose
complexity grows non-linearly with the problem size. For this reason, regression results should be
interpreted as an empirical approximation of observed runtime behaviour, rather than as a substitute
for theoretical complexity analysis.

\subsection{Regularized Linear Models}
\label{subsec:regularized_linear_models}
\hspace{\parindent}
When a regression model includes several explanatory variables, the fitted coefficients may become
sensitive to noise, correlated features, or limited training data. This can lead to overfitting, where
the model describes the training observations well but generalizes poorly to new data
\cite{theodoridis2020machine,hastie2009elements}. Regularization addresses this issue by adding a penalty
term to the regression objective, discouraging excessively large coefficients and improving model
stability.

Ridge regression extends ordinary least squares by adding an \(L_2\) penalty to the coefficient
values. Its objective function is given by

\begin{equation}
    \min_{\beta}
    \left[
    \sum_{i=1}^{n}
    (y_i - \hat{y}_i)^2
    +
    \lambda
    \sum_{j=1}^{p}
    \beta_j^2
    \right],
\end{equation}

where \(\lambda \geq 0\) controls the strength of the regularization. Larger values of \(\lambda\)
force the coefficients to become smaller, reducing model variance but potentially increasing bias.

Lasso regression instead uses an \(L_1\) penalty:

\begin{equation}
    \min_{\beta}
    \left[
    \sum_{i=1}^{n}
    (y_i - \hat{y}_i)^2
    +
    \lambda
    \sum_{j=1}^{p}
    |\beta_j|
    \right].
\end{equation}

Unlike Ridge regression, the \(L_1\) penalty can shrink some coefficients exactly to zero. As a
result, Lasso can perform implicit feature selection, which may be useful when several candidate
features are available but only a subset contributes meaningfully to runtime prediction
\cite{theodoridis2020machine,hastie2009elements}.

In the context of heuristic runtime modelling, regularized regression can be useful when multiple
scenario descriptors are considered simultaneously. For example, graph size, number of robots,
density-related measures, neighbourhood parameters, and conflict estimates may be correlated.
Regularization helps reduce the instability caused by such correlations and can improve the
generalization of the model to planning instances not used during fitting.

Nevertheless, regularized models introduce an additional hyperparameter, \(\lambda\), which must
be selected empirically, commonly through validation data or cross-validation. Therefore, in this
work, regularized regression should be interpreted as a complementary empirical tool for runtime
analysis, while theoretical complexity remains the primary method for characterizing worst-case
algorithmic behaviour.

\subsection{Evaluation Metrics for Regression Models}
\label{subsec:regression_metrics}
\hspace{\parindent}
Regression models are typically evaluated by comparing predicted values with observed values on
data not used to fit the model \cite{theodoridis2020machine}. In the context of runtime modelling, this
corresponds to comparing predicted execution times with measured execution times for independent
planning instances.

A commonly used metric is the Mean Absolute Error (MAE), defined as

\begin{equation}
    MAE
    =
    \frac{1}{n}
    \sum_{i=1}^{n}
    |y_i - \hat{y}_i|.
\end{equation}

MAE is easy to interpret because it is expressed in the same unit as the target variable. For runtime
prediction, this means that MAE directly represents the average absolute prediction error in units
of time.

Another common metric is the Root Mean Squared Error (RMSE), defined as

\begin{equation}
    RMSE
    =
    \sqrt{
    \frac{1}{n}
    \sum_{i=1}^{n}
    (y_i - \hat{y}_i)^2
    }.
\end{equation}

Compared with MAE, RMSE penalizes larger errors more strongly due to the squared term. This
makes it useful when large runtime prediction errors are particularly undesirable.

The coefficient of determination, usually denoted by \(R^2\), measures the proportion of variance
in the target variable explained by the model. It is defined as

\begin{equation}
    R^2
    =
    1
    -
    \frac{
    \sum_{i=1}^{n}
    (y_i - \hat{y}_i)^2
    }{
    \sum_{i=1}^{n}
    (y_i - \bar{y})^2
    },
\end{equation}

where \(\bar{y}\) is the mean of the observed target values. A value of \(R^2\) close to one indicates
that the model explains a large fraction of the observed variability. However, \(R^2\) should not be
interpreted alone, since a high value does not necessarily imply accurate predictions for all planning
instances.

For this reason, runtime models should be evaluated using a combination of error-based metrics,
such as MAE and RMSE, and goodness-of-fit measures, such as \(R^2\). This provides a more
complete view of both prediction accuracy and explanatory power.